\clearpage\fancyhead[L]{\small JUNO.}\hypertarget{native-377-lower}{}\label{native-377-lower}
\begin{center}\textbf{JUNO.}\par\vspace{4pt}\hrule\vspace{4pt}\fontsize{10}{12}\selectfont Göttingische gelehrte Anzeigen. Stück 204. Seite 2033\;..\;2034. 1813. December 23.\par\vspace{4pt}\hrule\vspace{4pt}\end{center}


\par\indent Die Jahrszeit, in welcher in diesem Jahre die \emph{Juno} in der \emph{Georgsharfe} mit der Sonne in Opposition kam, pflegt bei uns den astronomischen Beobachtungen nicht günstig zu sein: nur eine einzige Beobachtung dieses Planeten konnte Prof. \textsc{Gauss} um diese Zeit erhalten:
\begin{center}\fontsize{10}{12}\selectfont\setlength{\tabcolsep}{4pt}\renewcommand{\arraystretch}{1.15}\begin{tabular}{l|c|c|c}
1813 & Mittlere Zeit & Gerade Aufsteig. & Abweichung\\\hline
Novemb. 19 & $9^{\mathrm u}\;46^{\mathrm m}\;46''$ & $60^\circ\;30'\;35''0$ & $3^\circ\;5'\;9''3\ \text{Südl.}$\\
\end{tabular}\end{center}
\par\indent Obgleich der Planet diesmal zugleich sehr nahe bei seinem Perihelium, und also in Rücksicht seiner Erleuchtung in der allervortheilhaftesten Lage war, hatte er doch kaum die 8. Grösse.

\Needspace{14\baselineskip}\par\indent Da auswärtige Beobachtungen gänzlich fehlen, und gegenwärtig von entfernten Sternwarten, falls der Planet noch sonst wo beobachtet sein sollte, keine erwartet werden können, so musste sich Prof. \textsc{Gauss} entschliessen, die Berechnung der Opposition auf obige einzelne Beobachtungen zu gründen; das Resultat ist folgendes:
\begin{center}\leavevmode\SourceLeaf{378}\emph{Achte Opposition der Juno.}\end{center}
\begin{center}1813 November 19. $18^{\mathrm u}\;11^{\mathrm m}\;5^{\mathrm s}$ Mittlere Zeit in Göttingen.\end{center}
\begin{center}\fontsize{10}{12}\selectfont\setlength{\tabcolsep}{2pt}\renewcommand{\arraystretch}{1.15}\begin{tabular}{lcr}
Wahre Länge & $.\;.\;.\;.\;.\;.\;.\;.\;.\;.\;.\;.\;.$ & $57^\circ\;33'\;58''0$\\
Geocentrische Breite & $.\;.\;.\;.\;.\;.\;.\;.\;.\;.\;.$ & $23\;18\;46.8\ \text{Südl.}$\\
\end{tabular}\end{center}
\Needspace{20\baselineskip}\par\indent Zum Behuf einer für die nächste Sichtbarkeitsperiode des Planeten zu berechnenden Ephemeride, welche an einem andern Orte mitgetheilt werden wird, leitete Prof. \textsc{Gauss} aus der fünften, sechsten, siebenten und achten Opposition neue elliptische Elemente ab.
\begin{center}\fontsize{9}{11}\selectfont\setlength{\tabcolsep}{2pt}\renewcommand{\arraystretch}{1.15}\begin{tabular}{lcr}
\shortstack[l]{Epoche der mittlern Länge 1810\\im Meridian von Göttingen} &  & $95^\circ\;29'\;56''$\\
Tägliche tropische mittlere Bewegung & $.\;.\;.\;.\;.\;.\;.$ & $812''709$\\
Tropische Umlaufszeit & $.\;.\;.\;.\;.\;.\;.\;.\;.\;.\;.\;.\;.$ & $1594\ \text{Tage}\ 16\ \text{Stunden}$\\
Länge des Perihelium (1810) & $.\;.\;.\;.\;.\;.\;.\;.\;.\;.\;.$ & $53^\circ\;6'\;40''$\\
Länge des aufsteigenden Knoten (1810) & $.\;.\;.\;.\;.\;.\;.$ & $171\;7\;13$\\
Excentricitätswinkel & $.\;.\;.\;.\;.\;.\;.\;.\;.\;.\;.\;.\;.\;.$ & $14\;43\;8$\\
Neigung der Bahn & $.\;.\;.\;.\;.\;.\;.\;.\;.\;.\;.\;.\;.\;.$ & $13\;4\;12$\\
Logarithm. der halben grossen Axe & $.\;.\;.\;.\;.\;.\;.\;.$ & $0.4267670$\\
\end{tabular}\end{center}
\par\indent Die Vergleichung dieser Elemente mit denen, welche auf die vier \emph{ersten} Oppositionen gegründet waren (s. St. 136 dieser Anz. von 1808 [S. 306 d. B.]), zeigt den starken Einfluss der fremden Störungen auch bei diesem Planeten sehr sichtbar, besonders auffallend ist der Unterschied der mittlern Bewegung: die frühern Oppositionen gaben die Umlaufszeit um 3 Tage 4 Stunden kleiner. Ohne die Störungen hätte dieser Unterschied keine halbe Stunde betragen können.
